
\chapter{Experiments}

\section{Meaning representation}

\subsection{Data Collection}

The experimental evaluation for meaning representation requires structured data
from diverse domains. To obtain such data, \emph{quintd}~\cite{kasner2024beyond}
is used, a framework for collecting \emph{ad hoc} data-to-text datasets. 
Quintd selects examples using
a random seed and can also obtain up-to-date information from external data
sources.

The source collection used in the experiments is denoted by $S_0$. It contains
the five domains available in the corresponding quintd collection: GSMArena,
OpenWeather, Our World in Data (OWID), Wikidata, and ice hockey. The automatic
pipeline comparison uses all five domains. The human evaluation, fine-tuning,
and whole-system experiments use the four-domain subset consisting of GSMArena,
OpenWeather, OWID, and Wikidata. This distinction is important because ice
hockey is included in the automatic pipeline results but is absent from the
later $S_1$, $S_2$, $S_3$, and $S_4$ datasets.

The four experimental datasets are:

\begin{description}
  \item[\textbf{GSMArena}] Mobile phone specifications scraped from the GSMArena
  website. Each instance describes a single phone model with nested categories
  such as Network, Launch, Body, Display, Platform, Memory, Camera, and
  Battery, containing attribute--value pairs.

  \item[\textbf{OpenWeather}] Five-day weather forecasts obtained from the
  OpenWeather API. Each instance contains a time series of 40 three-hour
  intervals, with each interval recording temperature, humidity, wind speed,
  cloud coverage, and weather conditions.

  \item[\textbf{OWID}] Country-level time-series metrics from the Our World in
  Data project. Each instance represents a single metric, such as life
  expectancy, vaccination rates, or reproduction rate, for a specific country
  and contains yearly data points spanning several decades.

  \item[\textbf{Wikidata}] Knowledge-graph entities and their properties from
  Wikidata. Each instance describes an entity, such as a country, film, or
  organization, with a list of property--value pairs such as capital, cast
  member, or language used.
\end{description}

The fifth domain, \emph{ice hockey}, is therefore included only in the automatic
pipeline comparison. Its absence from the human evaluation, fine-tuning, and
whole-system experiments should not be interpreted as an omission from the
source collection.

All experimental datasets are provided as JSON files containing arrays of
structured records. A \emph{raw structured instance} is defined as a single
JSON object representing one self-contained data record from one of the four
experimental domains. Each instance contains nested fields with structured
data---time series, categorical values, numerical measurements, or hierarchical
attribute lists---that must be converted into semantic triples. Let
$D = \{d_1, \dots, d_N\}$ denote the set of $N$ such instances extracted from
the input JSON for a given domain.

\paragraph{Large language model.}
All LLM invocations across the three conversion methods use \emph{Gemma~4
31B IT}~\cite{team2024gemma,team2026gemma}, a 31-billion-parameter
instruction-tuned language model released by Google. The model is served via
\emph{vLLM}~\cite{kwon2023vllm}, a high-throughput inference engine that uses
\emph{PagedAttention} to manage key--value cache memory and enable batched
inference with dynamic sequence lengths.

\subsection{Experimental setup for pipelines}
\label{sec:experimental-setup}

The automatic pipeline comparison uses the $S_0$ source collection of 100 instances from each of the five
experimental domains. The reference-text generation stage is shared by the
three pipelines, while the text-to-triples stage is run separately for
Text-First Extraction with Normalization, Blueprint-Iterative Refinement, and
Evolving Catalog.

The generated outputs are evaluated using a human evaluation and an LLM as a judge. The judge is the
same Gemma~4 31B IT model, but it is used in a separate evaluation role rather
than as part of the extraction pipeline. Two transitions are assessed
independently:

\begin{enumerate}
    \item the conversion from the raw structured instance to the generated
    natural-language reference text; and
    \item the conversion from the reference text to the generated semantic
    triples.
\end{enumerate}

The two transitions use different pairs of criteria, with the complete prompts
provided in Appendix~\ref{app:pipeline-judge-prompts}. The data-to-text
transition is assessed for \emph{summary} and \emph{faithfulness}. Summary
measures whether the generated text is concise, coherent, non-redundant, and
preserves the core message of the structured source. Faithfulness measures
whether every textual claim is supported by that source, penalizing
hallucinations, contradictions, and material distortions. The text-to-triples
transition is assessed for \emph{additions} and \emph{omissions}. Additions
measures whether the triples introduce facts not supported by the reference
text, whereas omissions measures whether they leave out information whose
absence changes meaning, hides important context, or creates a misleading
imbalance. Every criterion is scored on an ordinal scale from 1 to 5, where 5
represents the best result.

The shared data-to-text evaluation requires $5 \times 100 \times 2 = 1{,}000$
judge calls. The three text-to-triples methods require
$5 \times 100 \times 3 \times 2 = 3{,}000$ further calls, for 4,000 calls in
total. For each domain and criterion, scores are averaged over the 100
instances. The resulting averages are used to compare the conversion methods.

The evaluation output also contains descriptive statistics. These statistics
are computed without additional LLM calls and describe the size and structure
of the generated outputs: the number of primitive JSON elements, reference
words, reference sentences, reference subsentences, extracted triples, and
unique predicates. The overall text and triples scores are calculated as the
mean of the two criteria belonging to the corresponding transition.

\begin{table}[!htbp]
\centering
\scriptsize
\caption{Automatic data-to-text evaluation shared by the three pipelines.}
\label{tab:pipeline-text-results}
\begin{tabular}{lrrrr}
\hline
\textbf{Domain} & \textbf{Summary} & \textbf{Faithfulness} & \textbf{Overall} & \textbf{JSON elements} \\
\hline
Ice hockey & 5.000 & 4.950 & 4.970 & 111.67 \\
GSMArena & 5.000 & 5.000 & 5.000 & 111.07 \\
OWID & 4.990 & 4.930 & 4.960 & 478.06 \\
OpenWeather & 5.000 & 5.000 & 5.000 & 902.92 \\
Wikidata & 4.990 & 4.990 & 4.990 & 13.18 \\
\hline
\textbf{Average} & 4.996 & 4.974 & 4.984 & 323.38 \\
\hline
\end{tabular}
\end{table}

\newcommand{\tripleheader}{\textbf{Domain} & \textbf{Additions} & \textbf{Omissions} & \textbf{Overall} & \textbf{Ref. words} & \textbf{Ref. sent.} & \textbf{Ref. subsent.} & \textbf{Triples} & \textbf{Predicates} \\}
\begin{table}[!htbp]
\centering
\scriptsize
\caption{Automatic text-to-triples evaluation for Blueprint-Iterative Refinement.}
\label{tab:pipeline-blueprint-triples-results}
\resizebox{\textwidth}{!}{%
\begin{tabular}{lrrrrrrrr}
\hline
\tripleheader
\hline
Ice hockey & 3.580 & 3.380 & 3.480 & 34.07 & 1.35 & 4.26 & 4.79 & 12 \\
GSMArena & 4.960 & 3.840 & 4.400 & 70.74 & 3.34 & 9.05 & 12.96 & 25 \\
OWID & 4.940 & 4.810 & 4.880 & 57.18 & 2.27 & 6.22 & 15.87 & 8 \\
OpenWeather & 4.880 & 3.840 & 4.360 & 66.21 & 2.82 & 6.53 & 5.46 & 18 \\
Wikidata & 4.910 & 4.410 & 4.660 & 20.34 & 1.17 & 2.61 & 4.15 & 47 \\
\hline
\textbf{Average} & 4.654 & 4.056 & 4.356 & 49.71 & 2.19 & 5.73 & 8.65 & 22.0 \\
\hline
\end{tabular}
}%
\end{table}

\begin{table}[!htbp]
\centering
\scriptsize
\caption{Automatic text-to-triples evaluation for Text-First Extraction with Normalization.}
\label{tab:pipeline-text-first-triples-results}
\resizebox{\textwidth}{!}{%
\begin{tabular}{lrrrrrrrr}
\hline
\tripleheader
\hline
Ice hockey & 5.000 & 4.970 & 4.990 & 32.54 & 1.32 & 4.21 & 6.77 & 106 \\
GSMArena & 5.000 & 5.000 & 5.000 & 70.34 & 3.26 & 9.30 & 13.61 & 159 \\
OWID & 5.000 & 4.990 & 5.000 & 58.23 & 2.28 & 6.31 & 8.09 & 281 \\
OpenWeather & 5.000 & 5.000 & 5.000 & 66.54 & 2.83 & 6.65 & 10.78 & 349 \\
Wikidata & 4.970 & 5.000 & 4.990 & 20.00 & 1.15 & 2.55 & 4.22 & 111 \\
\hline
\textbf{Average} & 4.994 & 4.992 & 4.996 & 49.53 & 2.17 & 5.80 & 8.69 & 201.2 \\
\hline
\end{tabular}
}%
\end{table}

\begin{table}[!htbp]
\centering
\scriptsize
\caption{Automatic text-to-triples evaluation for Evolving Catalog.}
\label{tab:pipeline-catalog-triples-results}
\resizebox{\textwidth}{!}{%
\begin{tabular}{lrrrrrrrr}
\hline
\tripleheader
\hline
Ice hockey & 4.990 & 3.920 & 4.460 & 31.88 & 1.27 & 4.10 & 5.14 & 4 \\
GSMArena & 4.970 & 4.810 & 4.890 & 69.48 & 3.19 & 9.11 & 12.87 & 25 \\
OWID & 4.880 & 3.850 & 4.370 & 56.55 & 2.22 & 5.99 & 3.40 & 139 \\
OpenWeather & 4.830 & 4.320 & 4.580 & 65.62 & 2.77 & 6.77 & 8.84 & 10 \\
Wikidata & 4.810 & 4.780 & 4.790 & 20.19 & 1.14 & 2.62 & 3.69 & 45 \\
\hline
\textbf{Average} & 4.896 & 4.336 & 4.618 & 48.74 & 2.12 & 5.72 & 6.79 & 44.6 \\
\hline
\end{tabular}
}%
\end{table}

\FloatBarrier
\subsection{Human evaluation for pipelines}

Human evaluation follows the same asymmetric rubric as automatic evaluation.
The shared data-to-text transition is assessed for summary and faithfulness,
while each text-to-triples pipeline is assessed for additions and omissions.
For each of the four domains, ten instances are assessed for every relevant
transition. Scores are assigned by the author of the thesis on the same 1--5
scale as automatic evaluation.

\begin{table}[!htbp]
\centering
\caption{Human evaluation of the shared data-to-text transition.}
\label{tab:human-text-results}
\begin{tabular}{lrrr}
\hline
\textbf{Domain} & \textbf{Summary} & \textbf{Faithfulness} & \textbf{Overall} \\
\hline
GSMArena & 4.10 & 3.60 & 3.85 \\
OWID & 4.70 & 4.80 & 4.75 \\
OpenWeather & 4.60 & 3.60 & 4.10 \\
Wikidata & 4.60 & 4.80 & 4.70 \\
\hline
\textbf{Average} & 4.50 & 4.20 & 4.35 \\
\hline
\end{tabular}
\end{table}

\newcommand{\humantripleheader}{\textbf{Domain} & \textbf{Additions} & \textbf{Omissions} & \textbf{Overall} \\}
\begin{table}[!htbp]
\centering
\caption{Human evaluation of Blueprint-Iterative Refinement.}
\label{tab:human-blueprint-results}
\begin{tabular}{lrrr}
\hline
\humantripleheader
\hline
GSMArena & 4.30 & 3.70 & 4.00 \\
OWID & 4.00 & 4.30 & 4.15 \\
OpenWeather & 5.00 & 3.00 & 4.00 \\
Wikidata & 4.80 & 4.30 & 4.55 \\
\hline
\textbf{Average} & 4.525 & 3.825 & 4.175 \\
\hline
\end{tabular}
\end{table}

\begin{table}[!htbp]
\centering
\caption{Human evaluation of Text-First Extraction with Normalization.}
\label{tab:human-text-first-results}
\begin{tabular}{lrrr}
\hline
\humantripleheader
\hline
GSMArena & 4.90 & 5.00 & 4.95 \\
OWID & 5.00 & 4.50 & 4.75 \\
OpenWeather & 4.90 & 5.00 & 4.95 \\
Wikidata & 5.00 & 5.00 & 5.00 \\
\hline
\textbf{Average} & 4.950 & 4.875 & 4.913 \\
\hline
\end{tabular}
\end{table}

\begin{table}[!htbp]
\centering
\caption{Human evaluation of Evolving Catalog.}
\label{tab:human-catalog-results}
\begin{tabular}{lrrr}
\hline
\humantripleheader
\hline
GSMArena & 4.90 & 4.90 & 4.90 \\
OWID & 5.00 & 2.20 & 3.60 \\
OpenWeather & 4.90 & 3.70 & 4.30 \\
Wikidata & 4.80 & 4.80 & 4.80 \\
\hline
\textbf{Average} & 4.900 & 3.900 & 4.400 \\
\hline
\end{tabular}
\end{table}

\FloatBarrier
\subsection{Correlation between LLM and human evaluation}

The agreement between the LLM judge and the human evaluator is computed from
paired 1--5 scores.

The \emph{LLM score mean} and \emph{Human score mean} columns report the
arithmetic mean of the aligned score vector used for that row. For an additions
or omissions row, each mean is calculated from the ten scores for that single
criterion. For a pooled row, it is calculated from 20 scores: additions and
omissions for each of the ten instances. For a per-instance-mean row, it is
calculated from ten instance-level averages. These columns describe the rating
level of each evaluator and are not correlation coefficients.

Four complementary statistics are reported. Pearson's product--moment
correlation measures linear association between the raw values, but assumes
equal distances between the ordinal score levels. Spearman's $\rho$ measures
whether the LLM and human rank instances similarly and is the primary
association measure for the ordinal 1--5 scale. Kendall's $\tau$-b measures
pairwise rank concordance while correcting for ties, making it useful for the
small samples and repeated ratings in this experiment. Quadratic weighted
Cohen's kappa measures chance-corrected ordinal agreement and penalizes a
two-point disagreement more strongly than a one-point disagreement
\cite{cohen1960agreement}. All four measures range from $-1$ to $1$; positive
values indicate agreement or positive association, zero indicates no
association or agreement beyond chance, and negative values indicate inverse
association.

The overall variants are calculated separately for the two evaluation tasks and
never mix data-to-text with text-to-triples. A \emph{pooled} text-to-triples
result concatenates the additions and omissions score of every matched instance
for each rater. With ten matched instances, this produces 20 paired criterion
scores: additions and omissions for each instance. It preserves
criterion-specific disagreement, for example when a judge agrees on additions
but not on omissions.

The complementary \emph{per-instance mean} result first averages additions and
omissions for each rater within an instance and then correlates the resulting
ten overall-quality pairs. It answers whether the LLM and human agree about
which complete instances are better or worse overall, while deliberately
discarding criterion-specific differences. For this mean variant, Pearson,
Spearman, and Kendall are computed from the floating-point means; quadratic
weighted kappa uses the means rounded to the nearest ordinal score.

Pearson, Spearman, and Kendall are undefined when either rater assigns the same
score to every observation. The status column records such constant-vector
cases. Means remain available for comparing score levels, but are not
correlation coefficients. In particular, a kappa of zero in a constant-vector
case indicates no measurable agreement beyond chance under the observed score
distributions; it does not establish that every individual rating disagreed.

\begin{table}[!htbp]
\centering
\scriptsize
\caption{Data-to-text correlation diagnostics ($N=10$ instances per row).}
\label{tab:judge-human-text-correlation}
\resizebox{\textwidth}{!}{%
\begin{tabular}{llrrrrrrl}
\hline
\textbf{Domain} & \textbf{Criterion} & \textbf{Pearson $r$} & \textbf{Spearman $\rho$} & \textbf{Kendall $\tau$} & \textbf{QW-$\kappa$} & \textbf{LLM score mean} & \textbf{Human score mean} & \textbf{Status} \\
\hline
GSMArena & Summary & -- & -- & -- & 0.000 & 5.00 & 4.10 & LLM constant \\
GSMArena & Faithfulness & -- & -- & -- & 0.000 & 5.00 & 3.60 & LLM constant \\
OWID & Summary & -- & -- & -- & 0.000 & 5.00 & 4.70 & LLM constant \\
OWID & Faithfulness & -- & -- & -- & 0.000 & 5.00 & 4.80 & LLM constant \\
OpenWeather & Summary & -- & -- & -- & 0.000 & 5.00 & 4.60 & LLM constant \\
OpenWeather & Faithfulness & -- & -- & -- & 0.000 & 5.00 & 3.60 & LLM constant \\
Wikidata & Summary & -- & -- & -- & 0.000 & 5.00 & 4.60 & LLM constant \\
Wikidata & Faithfulness & -- & -- & -- & 0.000 & 5.00 & 4.80 & LLM constant \\
\hline
\end{tabular}
}%
\end{table}

\begin{table}[!htbp]
\centering
\tiny
\caption{Text-to-triples correlation results by criterion and aggregation.}
\label{tab:judge-human-triples-correlation}
\resizebox{\textwidth}{!}{%
\begin{tabular}{lllrrrrrrrl}
\hline
\textbf{Pipeline} & \textbf{Domain} & \textbf{Criterion / aggregation} & \textbf{$N$} & \textbf{$r$} & \textbf{$\rho$} & \textbf{$\tau$} & \textbf{QW-$\kappa$} & \textbf{LLM score mean} & \textbf{Human score mean} & \textbf{Status} \\
\hline
Blueprint & GSMArena & Additions & 10 & -0.509 & -0.509 & -0.509 & -0.212 & 4.90 & 4.30 & ok \\
 &  & Omissions & 10 & 0.349 & 0.387 & 0.360 & 0.286 & 4.00 & 3.70 & ok \\
 &  & Pooled & 20 & 0.402 & 0.371 & 0.344 & 0.316 & 4.45 & 4.00 & ok \\
 &  & Per-instance mean & 10 & 0.000 & 0.000 & 0.000 & 0.000 & 4.45 & 4.00 & ok \\
 & OWID & Additions & 10 & -- & -- & -- & 0.000 & 5.00 & 4.00 & both constant, different \\
 &  & Omissions & 10 & 0.512 & 0.520 & 0.494 & 0.314 & 4.80 & 4.30 & ok \\
 &  & Pooled & 20 & 0.379 & 0.360 & 0.348 & 0.133 & 4.90 & 4.15 & ok \\
 &  & Per-instance mean & 10 & 0.512 & 0.520 & 0.494 & 0.000 & 4.90 & 4.15 & ok \\
 & OpenWeather & Additions & 10 & -- & -- & -- & 0.000 & 4.70 & 5.00 & Human constant \\
 &  & Omissions & 10 & 0.588 & 0.529 & 0.494 & 0.462 & 3.50 & 3.00 & ok \\
 &  & Pooled & 20 & 0.726 & 0.731 & 0.664 & 0.714 & 4.10 & 4.00 & ok \\
 &  & Per-instance mean & 10 & 0.633 & 0.536 & 0.501 & 0.000 & 4.10 & 4.00 & ok \\
 & Wikidata & Additions & 10 & 0.667 & 0.667 & 0.667 & 0.615 & 4.90 & 4.80 & ok \\
 &  & Omissions & 10 & 0.927 & 0.964 & 0.934 & 0.906 & 4.40 & 4.30 & ok \\
 &  & Pooled & 20 & 0.895 & 0.906 & 0.886 & 0.873 & 4.65 & 4.55 & ok \\
 &  & Per-instance mean & 10 & 0.907 & 0.868 & 0.841 & 0.800 & 4.65 & 4.55 & ok \\
\hline
Text-First & GSMArena & Additions & 10 & -- & -- & -- & 0.000 & 5.00 & 4.90 & LLM constant \\
 &  & Omissions & 10 & -- & -- & -- & 0.000 & 5.00 & 5.00 & both constant, same \\
 &  & Pooled & 20 & -- & -- & -- & 0.000 & 5.00 & 4.95 & LLM constant \\
 &  & Per-instance mean & 10 & -- & -- & -- & 0.000 & 5.00 & 4.95 & LLM constant \\
 & OWID & Additions & 10 & -- & -- & -- & 0.000 & 5.00 & 5.00 & both constant, same \\
 &  & Omissions & 10 & -- & -- & -- & 0.000 & 5.00 & 4.50 & LLM constant \\
 &  & Pooled & 20 & -- & -- & -- & 0.000 & 5.00 & 4.75 & LLM constant \\
 &  & Per-instance mean & 10 & -- & -- & -- & 0.000 & 5.00 & 4.75 & LLM constant \\
 & OpenWeather & Additions & 10 & -- & -- & -- & 0.000 & 5.00 & 4.90 & LLM constant \\
 &  & Omissions & 10 & -- & -- & -- & 0.000 & 5.00 & 5.00 & both constant, same \\
 &  & Pooled & 20 & -- & -- & -- & 0.000 & 5.00 & 4.95 & LLM constant \\
 &  & Per-instance mean & 10 & -- & -- & -- & 0.000 & 5.00 & 4.95 & LLM constant \\
 & Wikidata & Additions & 10 & -- & -- & -- & 0.000 & 5.00 & 5.00 & both constant, same \\
 &  & Omissions & 10 & -- & -- & -- & 0.000 & 5.00 & 5.00 & both constant, same \\
 &  & Pooled & 20 & -- & -- & -- & 0.000 & 5.00 & 5.00 & both constant, same \\
 &  & Per-instance mean & 10 & -- & -- & -- & 0.000 & 5.00 & 5.00 & both constant, same \\
\hline
Catalog & GSMArena & Additions & 10 & -- & -- & -- & 0.000 & 5.00 & 4.90 & LLM constant \\
 &  & Omissions & 10 & 0.667 & 0.667 & 0.667 & 0.615 & 4.80 & 4.90 & ok \\
 &  & Pooled & 20 & 0.444 & 0.444 & 0.444 & 0.444 & 4.90 & 4.90 & ok \\
 &  & Per-instance mean & 10 & 0.375 & 0.375 & 0.375 & 0.375 & 4.90 & 4.90 & ok \\
 & OWID & Additions & 10 & -- & -- & -- & 0.000 & 4.60 & 5.00 & Human constant \\
 &  & Omissions & 10 & 0.867 & 0.917 & 0.875 & 0.404 & 3.60 & 2.20 & ok \\
 &  & Pooled & 20 & 0.604 & 0.630 & 0.578 & 0.559 & 4.10 & 3.60 & ok \\
 &  & Per-instance mean & 10 & 0.627 & 0.795 & 0.707 & 0.043 & 4.10 & 3.60 & ok \\
 & OpenWeather & Additions & 10 & -- & -- & -- & 0.000 & 5.00 & 4.90 & LLM constant \\
 &  & Omissions & 10 & 0.753 & 0.646 & 0.631 & 0.478 & 4.40 & 3.70 & ok \\
 &  & Pooled & 20 & 0.782 & 0.647 & 0.632 & 0.630 & 4.70 & 4.30 & ok \\
 &  & Per-instance mean & 10 & 0.591 & 0.410 & 0.387 & 0.000 & 4.70 & 4.30 & ok \\
 & Wikidata & Additions & 10 & -- & -- & -- & 0.000 & 5.00 & 4.80 & LLM constant \\
 &  & Omissions & 10 & 0.667 & 0.667 & 0.667 & 0.615 & 4.80 & 4.80 & ok \\
 &  & Pooled & 20 & 0.459 & 0.459 & 0.459 & 0.444 & 4.90 & 4.80 & ok \\
 &  & Per-instance mean & 10 & 0.804 & 0.646 & 0.626 & 0.412 & 4.90 & 4.80 & ok \\
\hline
\end{tabular}
}%
\end{table}

\FloatBarrier
\subsection{Experimental setup for fine-tuning}

The fine-tuning experiments are conducted on the same four domains as the
pipeline comparison: GSMArena, OpenWeather, OWID, and Wikidata. The
training-and-validation collection is downloaded using seed $S_1$, while the
held-out test collection is downloaded independently using seed $S_2$.

The raw instances from the $S_1$ and $S_2$ collection are converted into reference text
and semantic triples using the Text-First Extraction with Normalization
pipeline. The aligned text--triple outputs are then used as the targets for a
domain-specific model. For each domain, the resulting examples are split into
80\% training data and 20\% validation data. The $S_2$ collection remains
separate and is used only for the final evaluation of the trained models. This
split prevents the test instances from influencing either model fitting or
training-configuration selection.

The fine-tuned model receives one structured instance and is trained to return
both a natural-language reference and its semantic triples in one JSON output.
The base model is Gemma~4 31B IT, and the adaptation is performed using QLoRA.
The base weights remain frozen, while low-rank adapters are trained in the
language-model attention and feed-forward projections. After training, the
adapters are merged into the base model before evaluation.

\subsection{Training Configuration}

The fine-tuning comparison contains five configurations. The standard baseline
uses three epochs. The \emph{baseline-epoch1} variant keeps the
baseline learning rate, adapter capacity, dropout, and weight decay, but trains
for one epoch. The remaining variants change the learning rate, adapter
capacity, or regularization in order to test their influence on the
domain-specialized models.

All training is performed on NVIDIA A100-SXM4 GPUs with 40~GB of VRAM, using
the PLGrid Athena cluster. The common effective batch size is 16, obtained
from a per-device batch size of 1 and 16 gradient-accumulation steps. The
optimizer is 8-bit paged AdamW, the learning-rate schedule is cosine, and the
compute precision is bfloat16. The base model is loaded using 4-bit NF4
quantization with double quantization.

\begin{table}[!htbp]
\centering
\small
\caption{QLoRA configurations used for fine-tuning.}
\label{tab:qlora-config}
\begin{tabular}{lrrrrrrr}
\hline
\textbf{Configuration} & \textbf{Epochs} & \textbf{LR} & \textbf{Warmup} & \textbf{$r$} & \textbf{$\alpha$} & \textbf{Dropout} & \textbf{Weight decay} \\
\hline
baseline & 3 & $10^{-4}$ & 0.03 & 16 & 32 & 0.05 & 0.00 \\
baseline-epoch1 & 1 & $10^{-4}$ & 0.03 & 16 & 32 & 0.05 & 0.00 \\
low\_lr & 5 & $3\times10^{-5}$ & 0.05 & 16 & 32 & 0.05 & 0.01 \\
higher\_capacity & 3 & $5\times10^{-5}$ & 0.05 & 32 & 64 & 0.05 & 0.01 \\
regularized\_capacity & 5 & $3\times10^{-5}$ & 0.10 & 32 & 64 & 0.10 & 0.01 \\
\hline
\end{tabular}
\end{table}

The LoRA configuration targets all attention and MLP projection matrices in
the language-model backbone (\texttt{q\_proj}, \texttt{k\_proj},
\texttt{v\_proj}, \texttt{o\_proj}, \texttt{gate\_proj}, \texttt{up\_proj},
and \texttt{down\_proj}). A positively scoped module pattern excludes the
vision and audio towers of the multimodal checkpoint. This selects 420 target
modules and approximately 25 million trainable parameters, which is less than
0.1\% of the base model.

The maximum sequence length is adjusted per domain to accommodate the varying
sizes of the raw structured instances while staying within the GPU memory
budget. Wikidata uses 8192 tokens, GSMArena and OWID use 4096 tokens, and
OpenWeather uses 6144 tokens. Instances exceeding the selected maximum are
truncated.

After training, the LoRA adapter weights are combined with the frozen base
model weights to produce a standalone model. This process, computes $W_{\mathrm{final}} = W_0 + \Delta W$ for each adapted layer.
The merged model can be served by the same vLLM infrastructure as the base
model and produces both a natural-language verbalization and the corresponding
semantic triples in a single decode pass.

\FloatBarrier
\subsection{Evaluation of fine-tuned models}

The fine-tuned models are evaluated on the held-out $S_2$ test sets for each
domain. The evaluation compares the generated text with the reference text
using BLEU and METEOR. The generated and reference triples are compared as
sets of subject--predicate--object tuples. Micro-averaged triple precision,
recall, and F1 are computed from the resulting true-positive, false-positive,
and false-negative counts. Predicate adherence measures the proportion of
generated predicates that belong to the predicate catalog used to construct
the targets.

The results are reported for the base model, which is the Gemma 4 before any fine-tuning and separately for the five
fine-tuning configurations.

\newcommand{\finetuneheader}{\textbf{Domain} & \textbf{Text BLEU} & \textbf{Text METEOR} & \textbf{Triple precision} & \textbf{Triple recall} & \textbf{Triple F1} & \textbf{Predicate adherence} \\}

\begin{table}[!htbp]
\centering
\small
\caption{Base-model results on the held-out $S_2$ sets.}
\label{tab:finetune-base-results}
\resizebox{\textwidth}{!}{%
\begin{tabular}{lrrrrrr}
\hline
\finetuneheader
\hline
GSMArena & 0.5621 & 0.7455 & 0.0515 & 0.0413 & 0.0458 & 0.3863 \\
OpenWeather & 0.2574 & 0.4931 & 0.0013 & 0.0010 & 0.0011 & 0.3771 \\
OWID & 0.3919 & 0.6170 & 0.0000 & 0.0000 & 0.0000 & 0.4181 \\
Wikidata & 0.7474 & 0.9032 & 0.0305 & 0.0409 & 0.0349 & 0.3280 \\
\hline
\textbf{Average} & 0.4897 & 0.6897 & 0.0208 & 0.0208 & 0.0205 & 0.3774 \\
\hline
\end{tabular}
}%
\end{table}

\begin{table}[!htbp]
\centering
\small
\caption{Fine-tuned model results for the baseline configuration.}
\label{tab:finetune-baseline-results}
\resizebox{\textwidth}{!}{%
\begin{tabular}{lrrrrrr}
\hline
\finetuneheader
\hline
GSMArena & 0.6245 & 0.8297 & 0.4234 & 0.4411 & 0.4321 & 0.7087 \\
OpenWeather & 0.2121 & 0.4463 & 0.0010 & 0.0005 & 0.0007 & 0.2002 \\
OWID & 0.3096 & 0.5757 & 0.0245 & 0.0227 & 0.0236 & 0.7148 \\
Wikidata & 0.7642 & 0.9087 & 0.5137 & 0.4495 & 0.4795 & 0.6635 \\
\hline
\textbf{Average} & 0.4776 & 0.6901 & 0.2407 & 0.2285 & 0.2340 & 0.5718 \\
\hline
\end{tabular}
}%
\end{table}

\begin{table}[!htbp]
\centering
\small
\caption{Fine-tuned model results for the baseline-epoch1 configuration.}
\label{tab:finetune-baseline-epoch1-results}
\resizebox{\textwidth}{!}{%
\begin{tabular}{lrrrrrr}
\hline
\finetuneheader
\hline
GSMArena & 0.5937 & 0.7969 & 0.2936 & 0.2942 & 0.2939 & 0.8389 \\
OpenWeather & 0.3175 & 0.5430 & 0.0013 & 0.0010 & 0.0011 & 0.3199 \\
OWID & 0.3745 & 0.6355 & 0.0090 & 0.0079 & 0.0084 & 0.6147 \\
Wikidata & 0.7136 & 0.8894 & 0.4103 & 0.1154 & 0.1801 & 0.7479 \\
\hline
\textbf{Average} & 0.4998 & 0.7162 & 0.1786 & 0.1046 & 0.1209 & 0.6304 \\
\hline
\end{tabular}
}%
\end{table}

\begin{table}[!htbp]
\centering
\small
\caption{Fine-tuned model results for the low learning-rate configuration.}
\label{tab:finetune-low-lr-results}
\resizebox{\textwidth}{!}{%
\begin{tabular}{lrrrrrr}
\hline
\finetuneheader
\hline
GSMArena & 0.6399 & 0.8257 & 0.3645 & 0.3615 & 0.3630 & 0.7522 \\
OpenWeather & 0.3068 & 0.5421 & 0.0016 & 0.0010 & 0.0012 & 0.3033 \\
OWID & 0.3877 & 0.6457 & 0.0088 & 0.0078 & 0.0083 & 0.6047 \\
Wikidata & 0.7752 & 0.9104 & 0.4363 & 0.4531 & 0.4446 & 0.7280 \\
\hline
\textbf{Average} & 0.5274 & 0.7310 & 0.2028 & 0.2059 & 0.2043 & 0.5971 \\
\hline
\end{tabular}
}%
\end{table}

\begin{table}[!htbp]
\centering
\small
\caption{Fine-tuned model results for the higher-capacity configuration.}
\label{tab:finetune-higher-capacity-results}
\resizebox{\textwidth}{!}{%
\begin{tabular}{lrrrrrr}
\hline
\finetuneheader
\hline
GSMArena & 0.6332 & 0.8316 & 0.4202 & 0.4310 & 0.4256 & 0.7306 \\
OpenWeather & 0.2448 & 0.4786 & 0.0027 & 0.0015 & 0.0019 & 0.1636 \\
OWID & 0.3654 & 0.6276 & 0.0256 & 0.0240 & 0.0248 & 0.6911 \\
Wikidata & 0.7513 & 0.9019 & 0.4757 & 0.4700 & 0.4728 & 0.6971 \\
\hline
\textbf{Average} & 0.4987 & 0.7099 & 0.2311 & 0.2316 & 0.2313 & 0.5706 \\
\hline
\end{tabular}
}%
\end{table}

\begin{table}[!htbp]
\centering
\small
\caption{Fine-tuned model results for the regularized-capacity configuration.}
\label{tab:finetune-regularized-capacity-results}
\resizebox{\textwidth}{!}{%
\begin{tabular}{lrrrrrr}
\hline
\finetuneheader
\hline
GSMArena & 0.6443 & 0.8297 & 0.3961 & 0.3953 & 0.3957 & 0.7373 \\
OpenWeather & 0.2957 & 0.5449 & 0.0025 & 0.0015 & 0.0018 & 0.2789 \\
OWID & 0.3911 & 0.6407 & 0.0191 & 0.0169 & 0.0179 & 0.6858 \\
Wikidata & 0.7555 & 0.8963 & 0.4912 & 0.4700 & 0.4803 & 0.6884 \\
\hline
\textbf{Average} & 0.5217 & 0.7279 & 0.2272 & 0.2209 & 0.2239 & 0.5976 \\
\hline
\end{tabular}
}%
\end{table}

\FloatBarrier
\section{Surface realization}

\subsection{WebNLG dataset and evaluation protocol}

Surface realization is evaluated on WebNLG, a benchmark in which RDF-style
semantic triples are paired with one or more natural-language lexicalizations.
The benchmark reader used in this work loads the English WebNLG release and
preserves both the triples and their reference texts. During evolution, the
training portion of WebNLG is used to provide fitness feedback. The final
programs are evaluated independently on the held-out test portion, preventing
the test references from being used during program search
\cite{gardent-etal-2017-webnlg}.

The evaluation is performed separately by WebNLG domain. The program receives
only the triples for each test entry. It must return one natural-language
string. The baseline systems are a rule-based generator by Lango et al.~~\cite{lango-dusek-2025-llm}, a fine-tuned BART
model, and a prompted language model. The prompt can be found in
Appendix~\ref{app:prompt-llm-judge}.

The final evaluation reports reference-based, judge-based, and descriptive
metrics. BLEU, BLEURT, and METEOR compare a generated text with the WebNLG
references. SENLEN measures the absolute difference between the number of
sentences in the generated text and the average number in the references. The
three judge columns contain scores from Gemma~4 31B IT, Qwen~3.6, and Themis.
The grammaticality, additions, and omissions judgments are performed with
Gemma~4 31B IT. The prompts used for these judgments are provided in
Appendix~\ref{app:prompt-llm-judge}. The final four columns are descriptive
statistics: average sentence count, word count, character count, and number
of input triples.

Higher values are preferable for BLEU, BLEURT, METEOR, judge scores, and
grammaticality. Lower values are preferable for additions, and
omissions; zero means that no addition or
omission was detected. Lower SENLEN values are preferable; zero indicates that
the period-delimited segment count exactly matches the average reference count.

\subsection{Evolution variants}

The initial all-domain run uses 1,000 iterations, Gemma~4 31B IT as the main
program-generation model, BLEU fitness, and one inspiration program in each
prompt. The later runs use 500 or 100 iterations and are restricted to three selected
domains.

The initial run was used to identify domains with different levels of
difficulty. WrittenWork was selected as the best-working domain, Building as
an intermediate domain, and Food as the most difficult domain. Restricting
subsequent variants to these three domains reduced the computational cost
of the experiments while retaining distinct task difficulties.

The following 500-iteration variants were then evaluated:

\begin{description}
    \item[BLEU without inspirations] BLEU fitness and no inspiration programs.
    \item[BLEU with inspirations] BLEU fitness and inspiration programs.
    \item[LLM judge without inspirations] LLM-as-a-judge fitness and no
    inspiration programs.
    \item[LLM judge with inspirations] LLM-as-a-judge fitness and inspiration
    programs.
    \item[Mixed-model LLM judge] LLM-as-a-judge fitness and inspirations, with
    Gemma~4 31B used for 80\% of calls and a 753B model for 20\% of calls.
    \item[Mixed-model LLM judge with four threads] The preceding mixed-model
    setup with four parallel evaluator threads.
\end{description}

Each run produces an executable surface-realization program. The program is
then evaluated on the held-out WebNLG test set rather than being judged only
by its online evolutionary fitness.

\subsection{Evolution results}

The following tables use the metric order shown in the evaluation protocol.
Average rows are calculated over the available rows for the
corresponding experiment.

\newcommand{\surfaceHeader}{\texttt{domain} & \texttt{BLEU} & \texttt{BLEURT} & \texttt{METEOR} & \texttt{SENLEN} & \texttt{Gemma} & \texttt{Qwen} & \texttt{Themis} & \texttt{gram.} & \texttt{add.} & \texttt{om.} & \texttt{sent.} & \texttt{words} & \texttt{chars} & \texttt{triples} \\\hline}
\newcommand{\surfacerow}[9]{#1 & #2 & #3 & #4 & #5 & #6 & #7 & #8 & #9 & \surfacetail}
\newcommand{\surfacetail}[6]{#1 & #2 & #3 & #4 & #5 & #6 \\\hline}
\newcommand{\surfacebaselinerow}[9]{#1 & #2 & #3 & #4 & #5 & #6 & #7 & #8 & #9 & \surfacebaselinetail}
\newcommand{\surfacebaselinetail}[7]{#1 & #2 & #3 & #4 & #5 & #6 & #7 \\\hline}

\begin{table}[!htbp]
\centering\tiny\setlength{\tabcolsep}{1pt}
\caption{Initial BLEU-based evolution variant with one inspiration.}
\label{tab:surface-base}
\begin{tabular}{l*{14}{r}}\hline
\surfaceHeader
\surfacerow{Airport}{0.4125}{0.2992}{0.7108}{0.4649}{0.9053}{0.7179}{0.8421}{0.3789}{0.0737}{0.0000}{0}{0}{0}{0}
\surfacerow{Artist}{0.3006}{0.1340}{0.6203}{0.2661}{0.8239}{0.6936}{0.7725}{0.5872}{0.0092}{0.2018}{0}{0}{0}{0}
\surfacerow{Astronaut}{0.5300}{0.2160}{0.7394}{0.4858}{0.9463}{0.7561}{0.8780}{0.6463}{0.0366}{0.0366}{0}{0}{0}{0}
\surfacerow{Athlete}{0.5632}{0.4583}{0.7828}{0.3000}{0.9360}{0.8280}{0.9240}{0.8200}{0.1000}{0.0000}{0}{0}{0}{0}
\surfacerow{Building}{0.4220}{0.0381}{0.6245}{0.8043}{0.9913}{0.8652}{0.9348}{0.5870}{0.0217}{0.0000}{0}{0}{0}{0}
\surfacerow{CelestialBody}{0.4962}{0.3329}{0.7172}{0.4898}{1.0000}{0.8286}{0.9510}{0.6122}{0.0000}{0.0000}{0}{0}{0}{0}
\surfacerow{City}{0.3463}{0.1817}{0.6536}{0.3976}{0.7831}{0.6843}{0.9157}{0.8072}{0.4578}{0.0000}{0}{0}{0}{0}
\surfacerow{ComicsCharacter}{0.2947}{-0.0221}{0.5241}{0.5500}{0.8000}{0.6067}{0.7733}{0.3000}{0.0000}{0.5667}{0}{0}{0}{0}
\surfacerow{Company}{0.4901}{0.2592}{0.7370}{0.3788}{0.9576}{0.7818}{0.9152}{0.5758}{0.0909}{0.0000}{0}{0}{0}{0}
\surfacerow{Food}{0.0764}{0.0417}{0.5173}{0.8587}{0.8652}{0.4913}{0.6348}{0.0000}{0.0435}{0.0000}{0}{0}{0}{0}
\surfacerow{MeanOfTransportation}{0.3523}{0.2908}{0.7313}{0.2126}{0.9241}{0.7517}{0.8966}{0.5517}{0.0345}{0.0345}{0}{0}{0}{0}
\surfacerow{Monument}{0.4950}{0.2289}{0.6997}{0.8877}{0.6043}{0.4261}{0.6652}{0.3913}{0.5870}{0.0000}{0}{0}{0}{0}
\surfacerow{Politician}{0.5059}{-0.0405}{0.7218}{0.7701}{1.0000}{0.8138}{0.9103}{0.1724}{0.0000}{0.0000}{0}{0}{0}{0}
\surfacerow{SportsTeam}{0.4638}{0.1271}{0.7756}{0.7992}{0.7818}{0.4773}{0.8136}{0.4545}{0.1136}{0.0000}{0}{0}{0}{0}
\surfacerow{University}{0.5473}{0.1314}{0.7261}{0.2926}{0.8733}{0.7889}{0.8644}{0.4444}{0.0889}{0.0111}{0}{0}{0}{0}
\surfacerow{WrittenWork}{0.5920}{0.2027}{0.7628}{0.3488}{0.9907}{0.8326}{0.9535}{0.5116}{0.0000}{0.0000}{0}{0}{0}{0}
\hline
\surfacerow{\textbf{Average}}{0.4305}{0.1800}{0.6903}{0.5192}{0.8864}{0.7090}{0.8528}{0.4900}{0.1036}{0.0532}{0}{0}{0}{0}
\end{tabular}
\end{table}

\begin{table}[!htbp]
\centering\tiny\setlength{\tabcolsep}{1pt}
\caption{BLEU-based evolution without inspirations, Gemma~4 31B, 500 iterations.}
\label{tab:surface-bleu-no-inspirations}
\begin{tabular}{l*{14}{r}}\hline
\surfaceHeader
\surfacerow{WrittenWork}{0.5565}{0.2162}{0.7654}{0.3488}{0.9814}{0.8930}{0.9581}{0.7442}{0.0000}{0.0000}{2.3488}{19.8140}{108.6744}{2.3488}
\surfacerow{Building}{0.4158}{0.0302}{0.6181}{0.8043}{0.9435}{0.7913}{0.8957}{0.3478}{0.1739}{0.0000}{3.3478}{34.7391}{196.8478}{4.5652}
\surfacerow{Food}{0.0841}{0.0683}{0.5106}{0.8587}{0.8609}{0.5087}{0.6913}{0.0000}{0.0435}{0.0000}{1.0000}{26.7609}{156.7391}{4.5652}
\hline
\surfacerow{\textbf{Average}}{0.3521}{0.1049}{0.6314}{0.6706}{0.9286}{0.7310}{0.8484}{0.3640}{0.0725}{0.0000}{2.2322}{27.1047}{154.0871}{3.8264}
\end{tabular}
\end{table}

\begin{table}[!htbp]
\centering\tiny\setlength{\tabcolsep}{1pt}
\caption{BLEU-based evolution with inspirations, Gemma~4 31B, 500 iterations.}
\label{tab:surface-bleu-with-inspirations}
\begin{tabular}{l*{14}{r}}\hline
\surfaceHeader
\surfacerow{WrittenWork}{0.5579}{0.1871}{0.7657}{0.4186}{0.9442}{0.7581}{0.9395}{0.4884}{0.0465}{0.0000}{2.4186}{19.3953}{107.3488}{2.3488}
\surfacerow{Building}{0.4158}{0.0302}{0.6181}{0.8043}{0.9652}{0.8609}{0.8783}{0.3478}{0.1522}{0.0435}{3.3478}{34.7391}{196.8478}{4.5652}
\surfacerow{Food}{0.0841}{0.0683}{0.5106}{0.8587}{0.8478}{0.4913}{0.7043}{0.0000}{0.1304}{0.0000}{1.0000}{26.7609}{156.7391}{4.5652}
\hline
\surfacerow{\textbf{Average}}{0.3526}{0.0952}{0.6315}{0.6939}{0.9191}{0.7034}{0.8407}{0.2787}{0.1097}{0.0145}{2.2555}{26.9651}{153.6453}{3.8264}
\end{tabular}
\end{table}

\begin{table}[!htbp]
\centering\tiny\setlength{\tabcolsep}{1pt}
\caption{LLM-judge evolution without inspirations, Gemma~4 31B, 500 iterations.}
\label{tab:surface-llm-no-inspirations}
\begin{tabular}{l*{14}{r}}\hline
\surfaceHeader
\surfacerow{WrittenWork}{0.5192}{0.2417}{0.7355}{0.4186}{1.0000}{0.9116}{0.9907}{0.9535}{0.0000}{0.0000}{2.4186}{20.4419}{106.9302}{2.3488}
\surfacerow{Building}{0.4198}{0.0804}{0.6118}{0.8043}{1.0000}{0.8087}{0.9043}{0.5870}{0.0000}{0.0000}{3.3478}{37.9565}{203.5435}{4.5652}
\surfacerow{Food}{0.0904}{-0.1788}{0.5009}{0.8587}{1.0000}{0.6435}{0.8435}{0.9348}{0.0000}{0.0000}{1.0000}{26.1522}{153.6739}{4.5652}
\hline
\surfacerow{\textbf{Average}}{0.3431}{0.0478}{0.6161}{0.6939}{1.0000}{0.7879}{0.9128}{0.8251}{0.0000}{0.0000}{2.2555}{28.1835}{154.7159}{3.8264}
\end{tabular}
\end{table}

\begin{table}[!htbp]
\centering\tiny\setlength{\tabcolsep}{1pt}
\caption{LLM-judge evolution with inspirations, Gemma~4 31B, 500 iterations.}
\label{tab:surface-llm-with-inspirations}
\begin{tabular}{l*{14}{r}}\hline
\surfaceHeader
\surfacerow{WrittenWork}{0.5405}{0.2194}{0.7667}{0.3488}{1.0000}{0.9511}{0.9904}{0.8733}{0.0000}{0.0000}{2.3488}{19.8140}{108.6744}{2.3488}
\surfacerow{Building}{0.4168}{0.0214}{0.6231}{0.8043}{1.0000}{0.8783}{0.9565}{0.5870}{0.0000}{0.0000}{3.3478}{34.7391}{196.8478}{4.5652}
\surfacerow{Food}{0.0671}{-0.2275}{0.4968}{0.8587}{0.7391}{0.4174}{0.5739}{0.0435}{0.0435}{0.0000}{1.0000}{27.2609}{157.1739}{4.5652}
\hline
\surfacerow{\textbf{Average}}{0.3414}{0.0044}{0.6289}{0.6706}{0.9130}{0.7489}{0.8403}{0.5013}{0.0145}{0.0000}{2.2322}{27.2713}{154.2321}{3.8264}
\end{tabular}
\end{table}

\begin{table}[!htbp]
\centering\tiny\setlength{\tabcolsep}{1pt}
\caption{Mixed-model LLM-judge evolution with inspirations, 100 iterations.}
\label{tab:surface-llm-mixed}
\begin{tabular}{l*{14}{r}}\hline
\surfaceHeader
\surfacerow{WrittenWork}{0.5499}{0.2371}{0.7615}{0.3488}{1.0000}{0.9721}{0.9953}{0.8837}{0.0000}{0.0000}{2.3488}{20.3721}{110.7907}{2.3488}
\surfacerow{Building}{0.3665}{0.0739}{0.6052}{0.8043}{1.0000}{0.8000}{0.9478}{0.8261}{0.0000}{0.0000}{3.3478}{38.5652}{212.0870}{4.5652}
\surfacerow{Food}{0.0938}{-0.0887}{0.4803}{0.8587}{1.0000}{0.5783}{0.8696}{0.3913}{0.0000}{0.0000}{1.0000}{28.2826}{170.1087}{4.5652}
\hline
\surfacerow{\textbf{Average}}{0.3367}{0.0741}{0.6156}{0.6706}{1.0000}{0.7835}{0.9376}{0.7004}{0.0000}{0.0000}{2.2322}{29.0733}{164.3288}{3.8264}
\end{tabular}
\end{table}

\begin{table}[!htbp]
\centering\tiny\setlength{\tabcolsep}{1pt}
\caption{Mixed-model LLM-judge evolution with inspirations and four threads, 100 iterations.}
\label{tab:surface-llm-mixed-four-threads}
\begin{tabular}{l*{14}{r}}\hline
\surfaceHeader
\surfacerow{WrittenWork}{0.5348}{0.2359}{0.7386}{0.4186}{1.0000}{0.9349}{0.9953}{0.8837}{0.0000}{0.0000}{2.4186}{19.8605}{104.6744}{2.3488}
\surfacerow{Building}{0.3664}{0.0301}{0.6198}{0.8043}{1.0000}{0.7870}{0.9217}{0.9565}{0.0000}{0.0000}{3.3478}{37.9565}{208.8478}{4.5652}
\surfacerow{Food}{0.0721}{-0.1559}{0.4824}{0.8587}{1.0000}{0.7565}{0.9174}{0.6087}{0.0000}{0.0000}{1.0000}{28.8043}{166.3478}{4.5652}
\hline
\surfacerow{\textbf{Average}}{0.3244}{0.0367}{0.6136}{0.6939}{1.0000}{0.8261}{0.9448}{0.8163}{0.0000}{0.0000}{2.2555}{28.8738}{159.9567}{3.8264}
\end{tabular}
\end{table}

\subsection{Configured experiments with unavailable results}

The following configurations were prepared after the runs reported above. They
use no inspiration programs and Gemma~4 31B IT as the LLM judge. Their result
artifacts were not available when this thesis was prepared. Each table is kept
as an explicit zero-filled placeholder so that the experimental matrix and its
intended reporting schema are documented; these zeros are not measured scores.

\newcommand{\surfaceplaceholderbody}{%
\surfacerow{WrittenWork}{0}{0}{0}{0}{0}{0}{0}{0}{0}{0}{0}{0}{0}{0}
\surfacerow{Building}{0}{0}{0}{0}{0}{0}{0}{0}{0}{0}{0}{0}{0}{0}
\surfacerow{Food}{0}{0}{0}{0}{0}{0}{0}{0}{0}{0}{0}{0}{0}{0}
\hline
\surfacerow{\textbf{Average}}{0}{0}{0}{0}{0}{0}{0}{0}{0}{0}{0}{0}{0}{0}}

\begin{table}[!htbp]
\centering\tiny\setlength{\tabcolsep}{1pt}
\caption{Configured 100-iteration mixed Gemma run: 80\% standard generation, 20\% high-reasoning generation, four evaluator threads (\texttt{6\_2gemmas\_thinking\_4thread\_pp}). Results unavailable.}
\label{tab:surface-config-6-two-gemma-four-thread}
\begin{tabular}{l*{14}{r}}\hline
\surfaceHeader
\surfaceplaceholderbody
\end{tabular}
\end{table}

\begin{table}[!htbp]
\centering\tiny\setlength{\tabcolsep}{1pt}
\caption{Configured 100-iteration mixed Gemma run: 80\% standard generation, 20\% high-reasoning generation, one evaluator thread (\texttt{6\_2gemmas\_thinking\_1thread\_pp}). Results unavailable.}
\label{tab:surface-config-6-two-gemma-one-thread-pp}
\begin{tabular}{l*{14}{r}}\hline
\surfaceHeader
\surfaceplaceholderbody
\end{tabular}
\end{table}

\begin{table}[!htbp]
\centering\tiny\setlength{\tabcolsep}{1pt}
\caption{Configured PLGrid 100-iteration mixed Gemma run: 80\% standard generation, 20\% high-reasoning generation, one evaluator thread (\texttt{6\_2gemmas\_thinking\_1thread\_plgrid}). Results unavailable.}
\label{tab:surface-config-6-two-gemma-one-thread-plgrid}
\begin{tabular}{l*{14}{r}}\hline
\surfaceHeader
\surfaceplaceholderbody
\end{tabular}
\end{table}

\begin{table}[!htbp]
\centering\tiny\setlength{\tabcolsep}{1pt}
\caption{Configured 500-iteration high-reasoning Gemma~4 31B run with four evaluator threads (\texttt{6\_1gemma\_thinking\_4thread\_500iter\_pp}). Results unavailable.}
\label{tab:surface-config-6-one-gemma-thinking}
\begin{tabular}{l*{14}{r}}\hline
\surfaceHeader
\surfaceplaceholderbody
\end{tabular}
\end{table}

\begin{table}[!htbp]
\centering\tiny\setlength{\tabcolsep}{1pt}
\caption{Configured 500-iteration Gemma~4 E4B run with one evaluator thread (\texttt{7\_gemma\_1thread\_500iter\_pp}). Results unavailable.}
\label{tab:surface-config-7-gemma}
\begin{tabular}{l*{14}{r}}\hline
\surfaceHeader
\surfaceplaceholderbody
\end{tabular}
\end{table}

\begin{table}[!htbp]
\centering\tiny\setlength{\tabcolsep}{1pt}
\caption{Configured 500-iteration Llama~3.3 70B Instruct AWQ INT4 run with one evaluator thread (\texttt{7\_llama\_1thread\_500iter\_pp}). Results unavailable.}
\label{tab:surface-config-7-llama}
\begin{tabular}{l*{14}{r}}\hline
\surfaceHeader
\surfaceplaceholderbody
\end{tabular}
\end{table}

\begin{table}[!htbp]
\centering\tiny\setlength{\tabcolsep}{1pt}
\caption{Configured 500-iteration Qwen3 32B AWQ run with one evaluator thread (\texttt{7\_qwen\_1thread\_500iter\_pp}). Results unavailable.}
\label{tab:surface-config-7-qwen}
\begin{tabular}{l*{14}{r}}\hline
\surfaceHeader
\surfaceplaceholderbody
\end{tabular}
\end{table}

\begin{table}[!htbp]
\centering\tiny\setlength{\tabcolsep}{1pt}
\caption{Configured 500-iteration standard Gemma~4 31B run with one evaluator thread (\texttt{8\_1gemma\_noref\_1thread\_500\_iter\_pp}). Results unavailable.}
\label{tab:surface-config-8-gemma-no-reasoning}
\begin{tabular}{l*{14}{r}}\hline
\surfaceHeader
\surfaceplaceholderbody
\end{tabular}
\end{table}

\FloatBarrier
\subsection{Baseline systems}

The baselines are a rule-based generator, a fine-tuned BART model, and an LLM
prompt. The average is calculated over the available domain results for each
baseline.

\begin{table}[!htbp]
\centering\tiny\setlength{\tabcolsep}{1pt}
\caption{Baseline scores across WebNLG domains.}
\label{tab:surface-baselines}
\resizebox{\textwidth}{!}{%
\begin{tabular}{ll*{14}{r}}\hline
\texttt{method} & \texttt{domain} & \texttt{BLEU} & \texttt{BLEURT} & \texttt{METEOR} & \texttt{SENLEN} & \texttt{Gemma} & \texttt{Qwen} & \texttt{Themis} & \texttt{gram.} & \texttt{add.} & \texttt{om.} & \texttt{sent.} & \texttt{words} & \texttt{chars} & \texttt{triples} \\\hline
\surfacebaselinerow{Rule-based generator}{all domains}{0.1702}{-0.3874}{0.5205}{1.0175}{0.4108}{0.4811}{0.7552}{0.1783}{0.0889}{0.0742}{2.5414}{21.7471}{146.6654}{3.1946}
\surfacebaselinerow{Fine-tuned BART}{all domains}{0.5077}{0.2599}{0.7163}{0.5157}{0.6366}{0.5933}{0.7899}{0.8366}{0.3142}{0.1863}{2.1679}{138.9440}{23.7167}{3.5205}
\surfacebaselinerow{LLM prompt}{all domains}{0.5019}{0.3443}{0.7246}{0.5274}{0.9954}{0.9743}{0.9847}{0.9861}{0.0045}{0.0022}{1.7307}{132.0601}{21.5757}{3.5205}
\hline
\end{tabular}
}%
\end{table}

\FloatBarrier
\subsection{Three-domain comparison}

The final comparison reports averages over WrittenWork, Building, and Food for
the baselines and all six reported three-domain OpenEvolve variants. The
all-domain initial variant is omitted because descriptive statistics were not
reported for that run.

\begin{table}[!htbp]
\centering\tiny\setlength{\tabcolsep}{1pt}
\caption{Average scores on WrittenWork, Building, and Food.}
\label{tab:surface-final-comparison}
\resizebox{\textwidth}{!}{%
\begin{tabular}{l*{14}{r}}\hline
\texttt{method} & \texttt{BLEU} & \texttt{BLEURT} & \texttt{METEOR} & \texttt{SENLEN} & \texttt{Gemma} & \texttt{Qwen} & \texttt{Themis} & \texttt{gram.} & \texttt{add.} & \texttt{om.} & \texttt{sent.} & \texttt{words} & \texttt{chars} & \texttt{triples} \\\hline
\surfacerow{Rule-based generator}{0.2063}{-0.4489}{0.5498}{1.9138}{0.2316}{0.3529}{0.6380}{0.0694}{0.4291}{0.1836}{4.3887}{30.8329}{199.2032}{3.3999}
\surfacerow{Fine-tuned BART}{0.5628}{0.2398}{0.7126}{0.7414}{0.6707}{0.5964}{0.8123}{0.7878}{0.2782}{0.1122}{2.3025}{147.8967}{27.3637}{3.8264}
\surfacerow{LLM prompt}{0.4781}{0.2310}{0.6720}{0.6706}{0.9848}{0.9660}{0.9826}{0.9778}{0.0000}{0.0000}{2.2322}{141.2164}{25.2174}{3.8264}
\hline
\surfacerow{BLEU, no inspirations}{0.3521}{0.1049}{0.6314}{0.6706}{0.9286}{0.7310}{0.8484}{0.3640}{0.0725}{0.0000}{2.2322}{27.1047}{154.0871}{3.8264}
\surfacerow{BLEU, with inspirations}{0.3526}{0.0952}{0.6315}{0.6939}{0.9191}{0.7034}{0.8407}{0.2787}{0.1097}{0.0145}{2.2555}{26.9651}{153.6453}{3.8264}
\surfacerow{LLM judge, no inspirations}{0.3431}{0.0478}{0.6161}{0.6939}{1.0000}{0.7879}{0.9128}{0.8251}{0.0000}{0.0000}{2.2555}{28.1835}{154.7159}{3.8264}
\surfacerow{LLM judge, with inspirations}{0.3414}{0.0044}{0.6289}{0.6706}{0.9130}{0.7489}{0.8403}{0.5013}{0.0145}{0.0000}{2.2322}{27.2713}{154.2321}{3.8264}
\surfacerow{Mixed-model LLM judge}{0.3367}{0.0741}{0.6156}{0.6706}{1.0000}{0.7835}{0.9376}{0.7004}{0.0000}{0.0000}{2.2322}{29.0733}{164.3288}{3.8264}
\surfacerow{Mixed-model LLM judge, 4 threads}{0.3244}{0.0367}{0.6136}{0.6939}{1.0000}{0.8261}{0.9448}{0.8163}{0.0000}{0.0000}{2.2555}{28.8738}{159.9567}{3.8264}
\end{tabular}
}%
\end{table}


\FloatBarrier
\section{Whole system evaluation}

The whole-system evaluation combines the meaning-representation and surface-
realization stages into one experimental pipeline. Unlike the preceding
experiments, which evaluate the two stages separately, this evaluation measures
whether errors introduced while converting raw structured data into triples
affect the final generated text.

\subsection{Data construction}

Two new collections are downloaded with quintd for each of the four domains:
GSMArena, OpenWeather, OWID, and Wikidata. The first collection is denoted by
$S_3$ and contains 1,000 raw structured instances per domain for training the
evolutionary programs. The second collection is denoted by $S_4$ and contains
another 1,000 raw structured instances per domain for final testing. The two
collections are downloaded independently so that no raw instance from the test
collection is available during evolution.

The raw instances in both collections are processed with the \emph{regularized capacity}
fine-tuned Gemma~4 model from the preceding fine-tuning experiments. For each
instance, the model produces a natural-language reference and the associated
normalized semantic triples in a single output. Thus, both $S_3$ and $S_4$
contain aligned pairs of the form
\[
    (T_i, r_i),
\]
where $T_i$ is the generated set of triples and $r_i$ is the corresponding
reference text. The model is selected before this evaluation and remains frozen
while the whole-system experiments are performed. The generated references are
therefore used as domain-specific evaluation targets, rather than as additional
inputs to the surface-realization programs.

\subsection{Evolution and evaluation protocol}

For each domain, the selected best-working OpenEvolve variant is run using the
$S_3$ collection. The evolved program receives only $T_i$ and must generate a
natural-language text. Its fitness is computed by comparing the generated text
with $r_i$ on the 1,000 training instances. A separate best program is retained
for each domain. The $S_4$ triples and reference texts are not used during
program generation, fitness calculation, or program selection.

After evolution, each domain-specific program is executed on all 1,000
instances in the corresponding $S_4$ test collection. The generated texts are
evaluated using the same protocol as the WebNLG surface-realization
experiments. BLEU, BLEURT, and METEOR compare the generated text with its
reference. SENLEN measures the difference between the generated and reference
sentence counts. Gemma~4, Qwen, and Themis provide judge scores, while
grammaticality, additions, and omissions are evaluated using the corresponding
judge prompts. Average sentence count, word count, character count, and number
of input triples are reported as descriptive statistics.

In this experiment, the reference triples are produced by the selected
fine-tuned model rather than annotated independently by humans. Consequently,
the evaluation measures the quality of the complete data-to-text pipeline with
respect to the distilled references. It does not separately estimate the
semantic precision or recall of the fine-tuned model; those properties are
reported in the fine-tuning experiments.

\begin{table}[!htbp]
\centering\tiny\setlength{\tabcolsep}{1pt}
\caption{Whole-system evaluation on the independent $S_4$ test collections.}
\label{tab:whole-system-results}
\begin{tabular}{l*{14}{r}}\hline
\surfaceHeader
\surfacerow{GSMArena}{0}{0}{0}{0}{0}{0}{0}{0}{0}{0}{0}{0}{0}{0}
\surfacerow{OpenWeather}{0}{0}{0}{0}{0}{0}{0}{0}{0}{0}{0}{0}{0}{0}
\surfacerow{OWID}{0}{0}{0}{0}{0}{0}{0}{0}{0}{0}{0}{0}{0}{0}
\surfacerow{Wikidata}{0}{0}{0}{0}{0}{0}{0}{0}{0}{0}{0}{0}{0}{0}
\hline
\surfacerow{\textbf{Average}}{0}{0}{0}{0}{0}{0}{0}{0}{0}{0}{0}{0}{0}{0}
\end{tabular}
\end{table}
